\chapter[AdaFace-Regional]{AdaFace-Regional: Proposta Arquitetônica}
\label{cap:adaface-regional}

Este capítulo descreve a contribuição arquitetônica do trabalho. O AdaFace-Regional foi projetado para verificar parentesco facial a partir de um backbone especializado em reconhecimento de identidade, mas adaptado à tarefa por meio de descongelamento parcial, comparação regional e regularização da ordem do par. O objetivo não é apenas trocar o extrator de características, e sim reorganizar o sinal global do AdaFace em uma representação comparativa voltada à tarefa de parentesco.

\section{Motivação}

A análise dos modelos avaliados no Capítulo~\ref{cap:Metodologia} indica que cada alternativa testa uma hipótese distinta, mas deixa uma limitação para o objetivo deste trabalho. O AdaFace-Cosseno usa um extrator facial forte, porém compara as faces apenas por similaridade global, sem adaptação supervisionada ao FIW e sem modelar quais partes do rosto participam da decisão. O ViT-FaCoR usa atenção cruzada entre patches, mas parte de um backbone visual geral e opera sobre muitos tokens locais, o que aumenta a liberdade do modelo para explorar regiões pouco informativas. O DINOv2-LoRA adapta uma representação auto-supervisionada ampla, mas não parte de um pré-treino facial especializado. O Modelo com Recuperação explicita exemplos positivos parecidos, mas depende da qualidade da galeria e da recuperação de suportes úteis.

A proposta AdaFace-Regional parte de uma hipótese intermediária: o sinal de reconhecimento facial do AdaFace já contém informação útil para parentesco e pode se beneficiar de uma reorganização para a tarefa. Para isso, o modelo mantém parte do backbone preservada, libera apenas os estágios finais para adaptação e compara regiões anatômicas do rosto de forma explícita. Além disso, como a verificação de parentesco é simétrica, a arquitetura busca tornar a decisão menos dependente da ordem em que as imagens aparecem no par.

\section{Visão geral da arquitetura}

A Figura~\ref{fig:adaface_regional_arquitetura} resume o fluxo do AdaFace-Regional. Cada face alinhada é processada pelo AdaFace IR-101 para produzir um embedding global e embeddings regionais. Os tokens regionais das duas faces passam por atenção cruzada bidirecional, gerando comparações por região. Um mecanismo de gating atribui pesos internos às regiões do par, e a cabeça binária final combina vetores globais, diferenças, produtos elemento a elemento, similaridades regionais e pesos do gate. Nas variantes simétricas, o mesmo par é avaliado nas ordens $(A,B)$ e $(B,A)$, e os logits são agregados antes da decisão.

\begin{figure}[htbp]
\centering
\fbox{%
\begin{minipage}{0.92\textwidth}
\centering
\small
\textbf{Par de entrada: faces alinhadas $A$ e $B$}\\[0.4em]
$\downarrow$\\[0.4em]
\begin{tabularx}{0.88\textwidth}{X c X}
AdaFace IR-101 em $A$ & & AdaFace IR-101 em $B$ \\
Embedding global $g_A$ & & Embedding global $g_B$ \\
Tokens regionais $r_A^1,\ldots,r_A^5$ & $\leftrightarrow$ & Tokens regionais $r_B^1,\ldots,r_B^5$ \\
\end{tabularx}\\[0.3em]
\begin{minipage}{0.84\textwidth}
\centering
Atenção cruzada entre regiões + similaridades $s_k$ + pesos $w_k$\\[0.2em]
Fusão comparativa: globais, diferenças, produtos, regiões e gate
\end{minipage}\\[0.4em]
$\downarrow$\\[0.4em]
\textbf{Logit kin/non-kin e cabeça auxiliar de relação}
\end{minipage}%
}
\caption{Fluxo arquitetônico do AdaFace-Regional.}
\label{fig:adaface_regional_arquitetura}
\end{figure}

\section{Backbone AdaFace}

O backbone do AdaFace-Regional é o AdaFace IR-101 \cite{kim2022adaface}, pré-treinado em WebFace4M para reconhecimento facial em larga escala. O modelo original produz embeddings de 512 dimensões, normalizados para comparação angular. Essa escolha é motivada pelo bom desempenho de embeddings faciais especializados em regimes de baixa falsa aceitação, observado na linha de base AdaFace-Cosseno.

No AdaFace-Regional, os três primeiros estágios do AdaFace permanecem congelados. Recebem gradiente apenas o estágio 4, formado pelos três blocos BasicBlockIR finais, e o \texttt{output\_layer}. Esse bloco descongelado soma cerca de 26 milhões de parâmetros treináveis. Com os módulos de regiões e o adaptador, as variantes completas totalizam 31,55 milhões de parâmetros treináveis de um total de 70,74 milhões, ou 44,6\%. A taxa de aprendizado do backbone é menor do que a da cabeça para reduzir o risco de descaracterizar representações de identidade já aprendidas.

\section{Tokens regionais}

Sobre a imagem alinhada em 224$\times$224 pixels, o AdaFace-Regional define cinco regiões fixas: face inteira, faixa superior dos olhos, região central do nariz, faixa média inferior da boca e faixa inferior da mandíbula. Cada região é recortada e redimensionada para 112$\times$112 antes de ser processada pelo AdaFace. Além do embedding global usado na fusão final, cada face gera cinco embeddings regionais; o token de face inteira é extraído como região e não reutiliza o embedding global.

Essa escolha introduz um viés estrutural simples. Em vez de permitir que o modelo compare livremente todos os patches da imagem, a comparação é orientada por partes faciais interpretáveis. O custo dessa decisão é uma dependência maior do alinhamento facial: se a detecção ou o recorte desloca as regiões, os tokens deixam de corresponder às partes esperadas.

\section{Atenção cruzada entre regiões}

Os cinco tokens de cada face passam por uma camada bidirecional de atenção cruzada entre regiões, com quatro cabeças. Para um par de faces, o módulo produz interações entre as regiões de $A$ e de $B$, formando uma matriz de atenção 5$\times$5. Essa matriz é muito menor do que os mapas de atenção entre patches dos modelos baseados em ViT, que operam tipicamente sobre centenas de posições.

A atenção regional permite que o modelo compare regiões correspondentes, como olhos com olhos, mas também teste interações entre regiões complementares. Assim, a arquitetura preserva a capacidade de interação entre faces sem abandonar a estrutura anatômica do problema.

\section{Mecanismo de gating e fusão}

Para cada região $k \in \{1, \ldots, K\}$, com $K = 5$, o modelo calcula uma similaridade regional $s_k \in \mathbb{R}$. Em paralelo, um módulo de gating sigmoide estima um peso $w_k \in [0,1]$ por região. Como o gate usa sigmoide e não softmax, mais de uma região pode receber peso alto ao mesmo tempo.

Na variante de referência, a cabeça final recebe um vetor de dimensão $4d + 2K + 1 = 2059$, com $d = 512$. Esse vetor concatena os embeddings globais $g_A$ e $g_B$, a diferença absoluta $|g_A - g_B|$, o produto elemento a elemento $g_A \odot g_B$, o vetor de similaridades regionais $s$, o vetor de pesos $w$ e o produto escalar $w^\top s$. A saída é um único logit, otimizado por entropia cruzada binária.

Uma variante posterior, chamada fusão apenas comparativa, remove os vetores globais crus $g_A$ e $g_B$ e mantém apenas termos que descrevem a comparação entre as faces. Com isso, a entrada da cabeça cai para 1035 dimensões. A intenção é verificar se a presença dos embeddings globais isolados ajuda a decisão ou se induz atalhos relacionados à aparência individual das faces.

\section{Cabeça auxiliar de relação}

Além da cabeça binária kin/non-kin, o AdaFace-Regional inclui uma cabeça auxiliar que prediz uma das onze relações positivas do FIW. Essa perda é aplicada apenas aos pares positivos, porque pares non-kin não têm tipo de relação anotado. O peso da perda auxiliar é 0,05 na configuração usada nas variantes completas.

A função dessa cabeça não é transformar o problema principal em classificação multiclasse. Ela atua como regularização supervisionada auxiliar, incentivando a representação interna a preservar diferenças anotadas no FIW entre relações como pai-filho, mãe-filha, irmãos e avós-netos. Essa informação é especialmente relevante porque as relações de segundo grau têm comportamento diferente das relações de primeiro grau.

\section{Passagem simétrica}

A verificação de parentesco é simétrica: se $A$ é parente de $B$, então $B$ é parente de $A$. Uma cabeça neural que recebe embeddings concatenados em ordem fixa, porém, pode aprender atalhos dependentes da posição da primeira e da segunda face no par.

Para reduzir esse problema, a variante simétrica processa cada par nas duas ordens, $(A,B)$ e $(B,A)$. Durante o treinamento, a perda binária é a média das duas direções. Durante a avaliação, o logit final é a média dos logits das duas permutações antes da aplicação do sigmoide:

\begin{equation}
z_{\text{sim}}(A,B) = \frac{z(A,B) + z(B,A)}{2}.
\label{eq:adaface_regional_simetrico}
\end{equation}

Essa passagem não impõe igualdade ponto a ponto entre $f(A,B)$ e $f(B,A)$. Ela reduz a dependência do escore final em relação à ordem de entrada, porque a decisão passa a considerar as duas permutações do mesmo par.

\section{Complexidade computacional}

O AdaFace-Regional adiciona custo principalmente em três pontos: extração dos embeddings regionais, atenção cruzada entre regiões e cabeça de fusão. O custo da atenção regional é pequeno, porque a matriz de interação tem apenas 5$\times$5 posições. O componente mais caro permanece sendo o próprio AdaFace IR-101, especialmente quando o estágio 4 e o \texttt{output\_layer} são descongelados.

Nas variantes completas, o modelo treina cerca de 31,55 milhões de parâmetros, contra 70,74 milhões de parâmetros totais. O número é menor do que um descongelamento integral do backbone, mas maior do que linhas de base que treinam apenas cabeças leves. Essa escolha reflete o compromisso central da proposta: adaptar o sinal facial especializado sem perder completamente as representações de identidade aprendidas no pré-treinamento.

\section{Variantes avaliadas}

Quatro variantes principais aparecem no ciclo experimental. A variante Base não usa a cabeça auxiliar de relação. Essa cabeça é introduzida no ciclo de desenvolvimento e mantida nas configurações finalistas Simétrica, Comparativa e com Negativos Difíceis.

\begin{table}[htbp]
\centering
\caption{Variantes avaliadas do AdaFace-Regional.}
\label{tab:adaface_regional_variantes}
\begin{tabularx}{\textwidth}{l X}
\hline
\textbf{Variante} & \textbf{Descrição} \\
\hline
Base & AdaFace-Regional com descongelamento parcial, tokens regionais, atenção cruzada, gate e cabeça binária em uma única ordem do par. \\
Simétrica & Mantém a cabeça auxiliar de relação, processa cada par nas duas ordens, $(A,B)$ e $(B,A)$, e agrega os logits para reduzir dependência da posição das imagens. \\
Comparativa & Mantém a cabeça auxiliar e a passagem simétrica, mas remove os embeddings globais crus da fusão final, preservando apenas termos comparativos entre as duas faces. \\
Negativos difíceis & Mantém a arquitetura comparativa e altera a amostragem de treino para incluir 30\% de negativos formados por famílias diferentes com os mesmos papéis familiares do par positivo de referência. \\
\hline
\end{tabularx}
\end{table}

As variantes não representam modelos independentes sem relação entre si. Elas formam uma sequência de decisões de projeto usada para testar se o ganho vem da adaptação parcial do AdaFace, da simetrização, da fusão comparativa ou de uma amostragem de negativos mais exigente. Os resultados dessas variantes são analisados no Capítulo~\ref{cap:resultados}.
